\subsection{Electrical subsystem}
\begin{figure}[ht]
    \centering
    \resizebox{0.5\textwidth}{!}{%
    \begin{circuitikz}
        \draw
        (0,0) to[V, v=$v_c(t)$] (0,4)
        to[R, l=$R_c$] (2.5,4)
        to[L, l=$L$] (5,4)
        to[R, l=$R_s$] (7.5,4)
        to[short, i=$i(t)$] (7.5,0)
        -- (0,0);
    \end{circuitikz}
    }
    \caption{Electrical model of the electromagnet circuit}
    \label{fig:rl_detailed}
\end{figure}

The electromagnet is modelled as a resistor-inductor circuit. If $v_c(t)$ is the voltage applied to the coil, the current dynamics are
\begin{equation}
    v_c(t) = R i(t) + L \frac{\dd i(t)}{\dd t}.
    \label{eq:electrical_rl}
\end{equation}

The resistor $R$ includes coil resistance $R_c$, sense resistance $R_s$.\\
\textcolor{red}{RESISTANCE CALIBRATION}
In addition the exact resistance value can be affected by variations and manufacturing tolerances.
Therefore the value is calibrated with the gain $R_{cal}$ obtained from the current measurement.
\[
R = R_{cal}(R_c + R_s)
\]

\textcolor{red}{INDUCTANCE CALIBRATION??}

When the amplifier gain is included, the coil voltage can be written as
\begin{equation}
    v_c(t) = K_g u(t),
\end{equation}
where $u(t)$ is the command voltage generated by the controller. Therefore,
\begin{equation}
    \dot{i}(t) = -\frac{R}{L}i(t) + \frac{K_g}{L}u(t).
    \label{eq:current_dynamics}
\end{equation}

The electrical plant is therefore a first-order transfer function:
\begin{equation}
G_i(s) = \frac{1}{Ls + R}
       = \frac{1}{R}\frac{\omega_i}{s + \omega_i},
\label{eq:electrical_plant_tf}
\end{equation}
where
\[
\omega_i = \frac{R}{L} \backsimeq 22,67~\mathrm{rad/s}
\]
is the characteristic frequency of the electrical subsystem, while the corresponding stable pole is
$p_i = -\omega_i$.
Then the time constant is $\tau_i = \frac{1}{\omega_i} \backsimeq 0.044~\mathrm{s}$.

\begin{figure}[H]
    \centering
    \includegraphics[width=0.9\textwidth]{../img/models/BODE_Gi.png}
    \caption{Bode diagram of $G_i$}
    \label{fig:bode_gi}
\end{figure}


\subsubsection{Current model validation}
\begin{figure}[H]
	\centering
	\includegraphics[width=0.9\textwidth]{../img/models/RL_model.png}
	\caption{Nonlinear RL model validation}
	\label{fig:rl_model_validation}
\end{figure}

From the comparison between simulated and measured responses, the final mean error is $0.00493\,\mathrm{A}$.
\begin{table}[H]
    \centering
    \begin{tabular}{cccc}
        \toprule
        Data & Rise time [s] & $\tau$ [s] & Pole [rad/s] \\
        \midrule
        Real & 0.082 & 0.0373 & -26.83 \\
        Simulation & 0.094 & 0.0427 & -23.40 \\
        \bottomrule
    \end{tabular}
    \caption{Estimated rise time, $\tau$, and pole from the step response.}
    \label{tab:current_tau_pole}
\end{table}

The pole obtained from $p_i=-\omega_i=-R/L$ is a nominal model value computed from the calibrated electrical
parameters. The pole in Table~\ref{tab:current_tau_pole} is instead an experimental estimate obtained
from the 10--90\% rise time, assuming a first-order step response:
\[
    \tau \simeq \frac{t_r}{\ln(9)}, \qquad p \simeq -\frac{1}{\tau}.
\]
The two values are therefore not expected to be exactly identical. The simulated response is close to
the nominal value, while the measured response can differ because of parameter uncertainty, amplifier
dynamics, sampling, noise.

